この研究では、人間が複数のsummaryを比較してどちらを好むかというdataを収集し、そのcomparisonからhuman preferenceを予測するmodelを学習する。そこから得たreward signalを使ってsummarization policyをreinforcement learningで調整する。重要なのは、training targetを既存referenceやROUGEのようなproxyだけに置かず、人間の比較判断そのものをoptimization loopへ入れたことにある。\autocite{src-8622fc1e22c2,src-e5838d162a66}


ただし、この結果を『general-purpose assistant alignmentが完成した』と読み替えてはいけない。対象はsummarizationであり、RedditのTL;DR dataを中心とした明確なtask boundaryがある。2020年の技術史として押さえるべきなのは、pretrained language modelの後段にhuman preferenceを使ったbehavior optimizationという操作点が実証的に置かれたことである。\autocite{src-8622fc1e22c2,src-e5838d162a66}


後年のRLHFという語から遡って見ると、preference data、reward model、reinforcement-learning policy optimizationという構成要素に連続性は見える。しかし2020年時点では、instruction-following assistant、large-scale deployment、安全性policy、system-level alignmentが一体化した後年のproduct stackはまだ存在しない。本号ではこの研究を『preference-feedback steeringの前段』として位置付け、後年の完成形を当年へ投影しない。\autocite{src-8622fc1e22c2,src-e5838d162a66}


Scaling LawsやGPT-3がpretrainingとin-context capabilityを、APIがaccessを、RAGがexternal knowledgeを扱ったのに対し、human-feedback summarizationはmodelの振る舞いをpretraining後にどう変えるかという別layerを示す。2020年には、生成AIを理解する単位がmodel architectureだけでは足りず、post-training objectiveまで含むsystemへ広がり始めていた。\autocite{src-8622fc1e22c2,src-e5838d162a66}


